\section{Servicio de red y conectividad interna}

\subsection{Contexto técnico}
La red interna del clúster implementa un modelo de enlaces punto a punto entre el nodo \texttt{master} y cada \texttt{worker}. Cada enlace usa una subred /28 independiente dentro de la superred \texttt{192.168.34.0/24}. Esta segmentación separa dominios de fallo y evita dependencias de switching L2 compartido para el tráfico de control interno. Sobre esa base se apoyan, en orden funcional, el montaje NFS por IP interna, el tráfico de control de Slurm y las validaciones de disponibilidad entre nodos.

La implementación no depende de archivos estáticos \path{/etc/sysconfig/network-scripts/ifcfg-*}. El repositorio opera sobre NetworkManager y usa \texttt{nmcli} como plano de control para crear, ajustar y activar conexiones persistentes. El motivo operativo es que la topología no se modela por plantilla fija de host, sino por inventario y variables declarativas; en ese contexto, escribir archivos \texttt{ifcfg} introduce dos riesgos: deriva de configuración cuando NM reescribe perfiles, y pérdida de idempotencia cuando el estado real del nodo no coincide con el archivo estático.

El estado esperado de red se expresa en variables globales. El rol \texttt{network\_internal} traduce ese estado en conexiones \texttt{int-*} y en entradas de \path{/etc/hosts}. El rol \texttt{cluster\_routing} añade enrutamiento persistente entre subredes internas y habilita el forward IPv4 en el \texttt{master}. El rol \texttt{firewall} aplica políticas de exposición mínima para Slurm sobre el CIDR interno.

Este servicio es de alto impacto porque toca conectividad base del clúster. Un valor incorrecto en interfaz, máscara o grupo de hosts no rompe sólo la red: propaga fallos a NFS, Slurm y validaciones posteriores. Por ese motivo el diseño incluye validación temprana con \texttt{assert}, lectura explícita de estado actual con \texttt{register} y tareas de consulta marcadas con \texttt{changed\_when: false}.

\subsection{Modelo declarativo de topología}
La topología se define en \path{group_vars/all/vars.yml} mediante un mapa \texttt{worker -> enlace}. Cada entrada describe ambos extremos del vínculo: interfaz e IP en \texttt{master} y en \texttt{worker}. Adicionalmente se declaran exclusiones para proteger interfaces de gestión o túneles superpuestos durante la limpieza de conexiones NetworkManager.

El siguiente fragmento corresponde al archivo de variables globales donde se declara la topología interna consumida por el rol.

\begin{filecode}{group_vars/all/vars.yml}
network_internal_keep_if: "eno1"
network_internal_exclude_ifaces:
  - "tailscale0"
network_internal_exclude_conn_regex: ".*tailscale.*"

network_internal_links:
  worker1:
    master_if: "enp3s0f1"
    master_ip: "192.168.34.1/28"
    worker_if: "enp3s0f1"
    worker_ip: "192.168.34.2/28"
  worker2:
    master_if: "enp3s0f2"
    master_ip: "192.168.34.17/28"
    worker_if: "enp3s0f2"
    worker_ip: "192.168.34.18/28"
  worker3:
    master_if: "enp3s0f3"
    master_ip: "192.168.34.33/28"
    worker_if: "enp3s0f3"
    worker_ip: "192.168.34.34/28"
  worker4:
    master_if: "enp3s0f0"
    master_ip: "192.168.34.49/28"
    worker_if: "enp3s0f0"
    worker_ip: "192.168.34.50/28"
\end{filecode}

A partir de ese mapa, el rol genera el conjunto de conexiones esperadas por nodo. En \texttt{master}, los nombres esperados son \texttt{int-<worker>}. En cada \texttt{worker} declarado, el nombre esperado es \texttt{int-master}. Esta convención evita ambigüedad al consultar o ajustar conexiones con \texttt{nmcli}.

La ausencia de una entrada para un nodo no detiene automáticamente el play si el nodo no es \texttt{master}; en ese caso la rama de configuración de enlace no se ejecuta. El efecto operativo es un nodo alcanzable por red de gestión pero no integrado a la malla interna del clúster.

\subsection{Variables del rol (tabla completa)}
\renewcommand{\arraystretch}{1.15}
\begin{longtable}{|>{\raggedright\arraybackslash}p{3.0cm}|>{\raggedright\arraybackslash}p{2.9cm}|>{\raggedright\arraybackslash}p{2.4cm}|>{\raggedright\arraybackslash}p{3.8cm}|>{\raggedright\arraybackslash}p{3.8cm}|}
\hline
\textbf{Variable} & \textbf{Origen / tipo} & \textbf{Obligatoria} & \textbf{Uso técnico en el rol} & \textbf{Efecto si está mal definida} \\
\hline
\endfirsthead
\hline
\textbf{Variable} & \textbf{Origen / tipo} & \textbf{Obligatoria} & \textbf{Uso técnico en el rol} & \textbf{Efecto si está mal definida} \\
\hline
\endhead
\texttt{network\_internal\_links} & Inventario, \texttt{mapping} & Sí & Define enlaces internos y genera conexiones esperadas; alimenta creación/ajuste de \texttt{int-*} y bloque en \path{/etc/hosts}. & El \texttt{assert} falla si no existe, no es \texttt{mapping} o está vacío. \\
\hline
\texttt{network\_internal\_links.<worker>.master\_if} & Inventario, \texttt{string} & Sí por enlace & Interfaz física del \texttt{master} usada en \texttt{nmcli con add}. & \texttt{nmcli} devuelve error de dispositivo inexistente y la tarea falla. \\
\hline
\texttt{network\_internal\_links.<worker>.master\_ip} & Inventario, CIDR & Sí por enlace & Dirección IPv4 del lado \texttt{master} y referencia para \path{/etc/hosts}. & IP/máscara inválida produce fallo en \texttt{nmcli}; máscara incorrecta rompe reachability y rutas posteriores. \\
\hline
\texttt{network\_internal\_links.<worker>.worker\_if} & Inventario, \texttt{string} & Sí por enlace & Interfaz física del \texttt{worker} en \texttt{int-master}. & Falla creación de conexión en worker por interfaz no encontrada. \\
\hline
\texttt{network\_internal\_links.<worker>.worker\_ip} & Inventario, CIDR & Sí por enlace & Dirección IPv4 del lado worker y entrada en \path{/etc/hosts}. & Error de validación en \texttt{nmcli} o conectividad parcial si la IP no corresponde a la subred declarada. \\
\hline
\texttt{network\_internal\_keep\_if} & Defaults/inventario, \texttt{string} & Sí & Interfaz de gestión que no debe tocarse en limpieza de conexiones. & El \texttt{assert} falla si es vacío; un valor incorrecto puede dejar sin protección la interfaz crítica. \\
\hline
\texttt{network\_internal\_exclude\_ifaces} & Defaults/inventario, lista & No & Excluye interfaces específicas del filtro de borrado. & Si falta, se usa lista vacía; si contiene interfaz interna por error, impide limpiar perfiles conflictivos. \\
\hline
\texttt{network\_internal\_exclude\_conn\_regex} & Defaults/inventario, \texttt{string} regex & No & Excluye conexiones por nombre (ej. Tailscale). & Regex demasiado amplia evita depuración; regex inválida puede abortar evaluación Jinja. \\
\hline
\texttt{network\_internal\_expected\_conns} & \texttt{set\_fact}, lista & Derivada & Lista canónica de conexiones internas por host para no borrarlas. & Si se calcula mal por inventario inconsistente, se pueden conservar perfiles obsoletos o no activar enlaces requeridos. \\
\hline
\texttt{nmcli\_conns} & \texttt{register} de shell & Derivada & Captura conexiones distintas de \texttt{keep\_if} para análisis de limpieza. & Si \texttt{nmcli} no está disponible o devuelve formato inesperado, \texttt{nmcli\_to\_delete} puede quedar vacío o incongruente. \\
\hline
\texttt{nmcli\_conn\_names} & \texttt{register} de command & Derivada & Detecta existencia de \texttt{int-master} o \texttt{int-<worker>}. & Con \texttt{failed\_when: false} no detiene play; puede ocultar falla previa y evitar creación/ajuste esperado. \\
\hline
\texttt{nmcli\_to\_delete} & \texttt{set\_fact}, lista & Derivada & Determina perfiles NM a eliminar antes de construir topología interna. & Cálculo vacío deja perfiles huérfanos; cálculo excesivo borra conexiones no objetivo. \\
\hline
\texttt{\_worker\_link} & \texttt{set\_fact}, \texttt{dict} & Derivada & Vista local del bloque \texttt{network\_internal\_links[inventory\_hostname]}. & Si el host no existe en el mapa no se define y no se crea enlace interno. \\
\hline
\texttt{worker\_link\_exists} / \texttt{master\_link\_exists} & \texttt{set\_fact}, boolean & Derivada & Controla si se ejecuta \texttt{nmcli con add}. & Valor incorrecto puede forzar recreaciones o impedir creación inicial. \\
\hline
\texttt{\_worker\_link\_needs\_mod} / \texttt{\_master\_link\_needs\_mod} & \texttt{set\_fact}, boolean & Derivada & Define si se ejecuta \texttt{nmcli con mod} según estado leído. & Si la lectura previa falla o está vacía, puede disparar modificaciones innecesarias. \\
\hline
\texttt{network\_internal\_hosts\_lines} & \texttt{set\_fact}, lista & Derivada & Construye bloque gestionado en \path{/etc/hosts}. & IPs sin normalizar o mapa incompleto generan resolución interna inconsistente. \\
\hline
\texttt{nmcli\_status} & \texttt{register} de command & Derivada & Reporte final de conexiones activas para depuración. & Al usar \texttt{failed\_when: false}, una falla puede quedar sólo como diagnóstico sin detener ejecución. \\
\hline
\end{longtable}

Las variables de integración que afectan el resultado final del servicio, aunque se consumen en roles posteriores, son \texttt{hpc\_internal\_subnets}, \texttt{hpc\_internal\_supernet}, \texttt{hpc\_router\_internal\_ifaces}, \texttt{slurm\_internal\_cidr}, \texttt{slurmctld\_port} y \texttt{slurmd\_port}. Si esas variables no son coherentes con \texttt{network\_internal\_links}, el enlace local puede quedar correcto pero el tráfico entre subredes o la exposición de puertos Slurm fallará.

\subsection{Secuencia operativa del rol \texttt{network\_internal}}
El rol aplica una secuencia de convergencia orientada a estado: valida insumos, descubre el estado real, calcula diferencias y sólo entonces crea o modifica conexiones.

\subsubsection{Validación temprana con \texttt{assert}}
La primera tarea impone un contrato mínimo de seguridad. La ejecución se corta antes de tocar NetworkManager si falta topología declarativa o si no se protegió la interfaz de gestión.

La validación de precondiciones se implementa en el archivo indicado mediante el siguiente bloque.

\begin{filecode}{roles/network_internal/tasks/main.yml}
- name: Red Interna | Assert variables definidas
  ansible.builtin.assert:
    that:
      - network_internal_links is defined
      - network_internal_links is mapping
      - network_internal_links | length > 0
      - network_internal_keep_if is defined
      - network_internal_keep_if is string
      - network_internal_keep_if | length > 0
    fail_msg: >-
      network_internal_links no está definido o está vacío, o network_internal_keep_if no está definido.
\end{filecode}

El \texttt{assert} protege contra ejecuciones peligrosas de dos tipos: limpieza de perfiles NM sin topología objetivo y operación sobre un host donde no se identificó una interfaz a preservar.

\subsubsection{Descubrimiento de estado con \texttt{register} y control de cambio}
Las tareas de lectura usan \texttt{register} para capturar salida y reutilizarla en condiciones posteriores. El patrón dominante es: consultar con \texttt{command}/\texttt{shell}, registrar resultado y marcar \texttt{changed\_when: false}. Así se separa explícitamente observación de mutación.

La lógica de descubrimiento de estado y registro de salidas de NetworkManager se define a continuación:

\begin{filecode}{roles/network_internal/tasks/main.yml}
- name: Red Interna | Obtener conexiones NM distintas de {{ network_internal_keep_if }}
  ansible.builtin.shell: |
    nmcli -t -f NAME,DEVICE con show | awk -F: '$2 != "{{ network_internal_keep_if }}" && $2 != "" {print $1}'
  register: nmcli_conns
  changed_when: false

- name: Red Interna | Listar conexiones NM (nombres)
  ansible.builtin.command: nmcli -g NAME con show
  register: nmcli_conn_names
  changed_when: false
  failed_when: false
\end{filecode}

\texttt{register} convierte salida textual en hechos de ejecución (\texttt{stdout}, \texttt{stdout\_lines}, \texttt{rc}) que luego alimentan \texttt{set\_fact} para cálculo de diferencias. \texttt{failed\_when: false} en tareas de inspección prioriza continuidad diagnóstica, pero debe interpretarse con cuidado porque evita cortar el play ante un fallo de consulta.

\subsubsection{Depuración de conexiones y uso de \texttt{nmcli}}
Con los datos registrados, el rol calcula \texttt{nmcli\_to\_delete} y elimina perfiles no permitidos antes de construir enlaces internos. Este paso previene colisiones de rutas o direcciones heredadas. A continuación, las operaciones de convergencia sobre cada enlace usan el ciclo \texttt{exists -> read -> compare -> mod -> up}.

En \texttt{master}, se itera cada entrada de \texttt{network\_internal\_links} mediante \path{master_link.yml}. En \texttt{worker}, se opera sobre \texttt{int-master}. El flujo de decisión es idéntico: sólo se crea si no existe y sólo se modifica si los parámetros actuales difieren de los esperados.

La implementación que materializa la convergencia por enlace en el master es la siguiente:

\begin{filecode}{roles/network_internal/tasks/master_link.yml}
- name: Red Interna | Crear conexion int-{{ link_item.key }}
  ansible.builtin.command: >-
    nmcli con add type ethernet ifname {{ link_item.value.master_if }} con-name int-{{ link_item.key }}
    ipv4.method manual ipv4.addresses {{ link_item.value.master_ip }} ipv6.method ignore
  when: not (master_link_exists | bool)
  register: _master_link_created

- name: Red Interna | Ajustar int-{{ link_item.key }}
  ansible.builtin.command: >-
    nmcli con mod int-{{ link_item.key }} ipv4.addresses {{ link_item.value.master_ip }} ipv4.method manual ipv6.method ignore
  when: _master_link_needs_mod
  register: _master_link_modified

- name: Red Interna | Activar int-{{ link_item.key }}
  ansible.builtin.command: nmcli con up int-{{ link_item.key }}
  changed_when: false
  when: _master_link_changed | default(false)
\end{filecode}

El uso de \texttt{nmcli} responde a dos requisitos del servicio: persistencia en perfiles de NetworkManager y aplicación inmediata sin reinicio global de red. La combinación \texttt{con add}, \texttt{con mod} y \texttt{con up} evita depender de scripts de arranque legacy y mantiene el estado alineado con el inventario.

\subsubsection{Actualización de resolución interna}
Tras configurar interfaces, el rol genera un bloque administrado en \path{/etc/hosts} con IPs internas sin máscara. Esto estabiliza resolución entre nodos para comandos operativos que usan nombres de inventario y reduce dependencia de DNS externo durante despliegue y validaciones internas.

El bloque se escribe con \texttt{blockinfile} y marcador propio, por lo que es idempotente y reversible. Si una IP en \texttt{network\_internal\_links} es incorrecta, la resolución local también lo será aunque la interfaz exista.

\subsection{Enrutamiento persistente (\texttt{cluster\_routing})}
El rol \texttt{cluster\_routing} resuelve el problema inter-subred: cada worker tiene conectividad directa sólo con su /28 local, por lo que necesita rutas explícitas hacia las demás subredes internas a través del \texttt{master} de su segmento.

En \texttt{master}, la tarea crítica habilita forwarding IPv4 persistente con \texttt{ansible.posix.sysctl} sobre \texttt{net.ipv4.ip\_forward=1}. Sin este paso, aunque existan rutas en workers, el tráfico no se reenviará entre enlaces internos.

En workers, la secuencia técnica es:
\texttt{(1)} detectar la interfaz cuyo \texttt{ansible\_facts[item].ipv4.address} coincide con \texttt{ansible\_host},
\texttt{(2)} resolver el nombre de conexión NM de esa interfaz,
\texttt{(3)} calcular subred local, gateway y destinos,
\texttt{(4)} comparar rutas deseadas con rutas actuales,
\texttt{(5)} agregar sólo diferencias y reactivar la conexión.

El cálculo de subred local, gateway y destinos se implementa en el siguiente fragmento:

\begin{filecode}{roles/cluster_routing/tasks/main.yml}
- name: "WORKERS | Calcular subnet local, gateway (master) y rutas destino"
  ansible.builtin.command: python3 -
  args:
    stdin: |
      import ipaddress, os, json
      ip = ipaddress.ip_address(os.environ["NODE_IP"])
      subnets = [ipaddress.ip_network(s) for s in json.loads(os.environ["SUBNETS_JSON"])]
      local = None
      for n in subnets:
        if ip in n:
          local = n
          break
      if local is None:
        raise SystemExit("NODE_IP no cae en hpc_internal_subnets")
      gw = ipaddress.ip_address(int(local.network_address) + 1)
      dest = [str(n) for n in subnets if n != local]
      print(str(local))
      print(str(gw))
      for n in dest:
        print(n)
  environment:
    NODE_IP: "{{ hostvars[inventory_hostname].ansible_host }}"
    SUBNETS_JSON: "{{ hpc_internal_subnets | to_json }}"
  register: _route_calc
  changed_when: false
\end{filecode}

El cálculo del gateway como \texttt{network\_address + 1} asume que el \texttt{master} ocupa esa IP en cada /28 según el contrato de \texttt{network\_internal\_links}. Si esa convención se rompe en variables, la ruta persistente apuntará a un next-hop inexistente.

La adición de rutas usa \texttt{+ipv4.routes} para no sobrescribir entradas previas y filtra por diferencia con \texttt{hpc\_current\_routes}. Esto reduce cambios innecesarios y mantiene idempotencia.

\subsection{Integración con firewall}
La integración de seguridad se divide en dos capas complementarias.

La primera capa la aplica \texttt{cluster\_routing}. En el \texttt{master}, pone interfaces internas en zona \texttt{trusted}. En workers, confía como origen la superred \texttt{hpc\_internal\_supernet} en la misma zona. Esta capa habilita tránsito interno base para forwarding y comunicación entre segmentos.

La segunda capa la aplica el rol \texttt{firewall}. Ahí se habilita \texttt{firewalld}, se permite SSH y se agregan \textit{rich rules} de Slurm condicionadas por grupo de inventario y CIDR interno.

La regla de control de acceso para el puerto de \texttt{slurmctld} se define en el archivo mostrado a continuación:

\begin{filecode}{roles/firewall/tasks/main.yml}
- name: Firewall | Slurm | Allow slurmctld (6817/tcp) from internal network (masters)
  ansible.posix.firewalld:
    zone: "{{ slurm_firewalld_zone }}"
    permanent: true
    immediate: true
    state: enabled
    rich_rule: >-
      rule family="ipv4"
      source address="{{ slurm_internal_cidr }}"
      port port="{{ slurmctld_port }}" protocol="tcp"
      accept
  when:
    - slurm_internal_cidr is defined
    - inventory_hostname in groups.get('hpc_master', [])
\end{filecode}

El uso de \texttt{notify} y \texttt{meta: flush\_handlers} fuerza recarga de \texttt{firewalld} antes de continuar, evitando ventanas donde las reglas estén declaradas pero aún no activas.

La coherencia entre \texttt{slurm\_internal\_cidr} y el direccionamiento efectivo de \texttt{network\_internal\_links} es obligatoria. Si no coinciden, los puertos Slurm están técnicamente abiertos pero no para el origen real de los nodos.

\subsection{Verificación operativa paso a paso}
La verificación se ejecuta por nodo y contrasta estado de interfaz, perfiles NM, tabla de rutas, reglas de firewall y puertos escuchando.

\subsubsection{Paso 1: inspección de interfaces (\texttt{ip a})}
Ejecutar en \texttt{master} y en cada \texttt{worker} para validar presencia de interfaces físicas internas y direcciones esperadas.

La inspección de interfaces y direccionamiento puede realizarse con los siguientes comandos:

\begin{terminalblock}
ip -br a
ip a show dev enp3s0f1
ip a show dev enp3s0f2
\end{terminalblock}

Si la dirección del enlace interno no aparece, revisar que la conexión \texttt{int-*} exista y esté activa en NetworkManager.

\subsubsection{Paso 2: estado de NetworkManager (\texttt{nmcli})}
Validar conexiones activas, su método IPv4/IPv6 y consistencia de nombres.

El estado operativo de conexiones de NetworkManager se comprueba ejecutando:

\begin{terminalblock}
nmcli -t -f NAME,DEVICE,TYPE,STATE con show --active
nmcli -g NAME con show
nmcli -g ipv4.addresses,ipv4.method,ipv6.method con show int-master
nmcli -g ipv4.addresses,ipv4.method,ipv6.method con show int-worker2
\end{terminalblock}

En \texttt{worker}, \texttt{int-master} debe quedar en \texttt{ipv4.method manual} e \texttt{ipv6.method ignore}. En \texttt{master}, cada \texttt{int-<worker>} debe mapear a la interfaz declarada en variables.

\subsubsection{Paso 3: validación de rutas (\texttt{ip route})}
Comprobar rutas locales y rutas a otras subredes internas.

La verificación de rutas efectivas se realiza mediante:

\begin{terminalblock}
ip route
ip route | grep 192.168.34
\end{terminalblock}

En workers se esperan rutas hacia subredes /28 remotas vía el gateway local (IP del \texttt{master} en esa subred). Si sólo aparece la ruta de enlace local, la fase \texttt{cluster\_routing} no convergió.

\subsubsection{Paso 4: validación de reglas de firewall (\texttt{firewall-cmd})}
Confirmar zonas activas, fuentes confiadas y reglas ricas de Slurm.

La política activa de \texttt{firewalld} puede inspeccionarse con:

\begin{terminalblock}
firewall-cmd --get-active-zones
firewall-cmd --zone=trusted --list-interfaces
firewall-cmd --zone=trusted --list-sources
firewall-cmd --zone=public --list-rich-rules
\end{terminalblock}

En \texttt{master}, deben observarse interfaces internas en \texttt{trusted}. En workers, debe observarse \texttt{192.168.34.0/24} como source confiada. Las \textit{rich rules} deben reflejar puertos 6817 o 6818 según rol de nodo.

\subsubsection{Paso 5: validación de sockets en escucha (\texttt{ss -lntp})}
Comprobar que los servicios dependientes de red interna están realmente escuchando.

La comprobación de puertos en escucha para servicios críticos se ejecuta con:

\begin{terminalblock}
ss -lntp | grep -E ':(6817|6818)\b'
ss -lntp | grep :22
\end{terminalblock}

\texttt{master} debe exponer \texttt{6817/tcp} para \texttt{slurmctld}. Los \texttt{workers} deben exponer \texttt{6818/tcp} para \texttt{slurmd}. Si el puerto no aparece, el problema no está en firewall sino en el servicio de aplicación.

\subsubsection{Ejecución remota de verificación desde nodo de control}
Cuando se necesite validación masiva sin entrar por SSH a cada nodo, se pueden lanzar comandos ad-hoc con Ansible.

La validación remota multipunto puede ejecutarse desde el nodo de control con:

\begin{terminalblock}
ansible -i inventario.ini all -b -m ansible.builtin.command -a "ip -br a"
ansible -i inventario.ini all -b -m ansible.builtin.command -a "nmcli -t -f NAME,DEVICE,TYPE,STATE con show --active"
ansible -i inventario.ini all -b -m ansible.builtin.command -a "ip route"
ansible -i inventario.ini all -b -m ansible.builtin.command -a "firewall-cmd --get-active-zones"
ansible -i inventario.ini all -b -m ansible.builtin.shell -a "ss -lntp | grep -E ':(6817|6818)\b' || true"
\end{terminalblock}

\subsection{Advertencias y escenarios de error}
\begin{criticalnote}
\textbf{Contrato de entrada incompleto.} Si \texttt{network\_internal\_links} está vacío o \texttt{network\_internal\_keep\_if} no tiene valor, el rol falla en la primera tarea por diseño. Forzar bypass de ese \texttt{assert} permite ejecutar limpieza y creación de conexiones sin topología validada, con riesgo de pérdida de conectividad de gestión.
\end{criticalnote}

\begin{criticalnote}
\textbf{Inconsistencia entre \texttt{ansible\_host} y red interna.} El cálculo de rutas en workers usa \texttt{ansible\_host} para ubicar la subred local dentro de \texttt{hpc\_internal\_subnets}. Si \texttt{ansible\_host} apunta a otra red, la tarea Python termina con \texttt{NODE\_IP no cae en hpc\_internal\_subnets}. El síntoma inmediato es ausencia de rutas persistentes a subredes remotas.
\end{criticalnote}

\begin{criticalnote}
\textbf{Interfaz física mal referenciada.} Un \texttt{master\_if} o \texttt{worker\_if} inexistente provoca error en \texttt{nmcli con add}. Incluso si el play continúa en otros hosts, el nodo afectado queda fuera de la malla interna y arrastra fallos de NFS/Slurm por hostname o por IP.
\end{criticalnote}

\begin{criticalnote}
\textbf{CIDR incompatible entre red y firewall.} Si \texttt{slurm\_internal\_cidr} o \texttt{hpc\_internal\_supernet} no corresponden a las IP reales de enlaces internos, las reglas se aplican en \texttt{firewalld} pero no autorizan los orígenes correctos. El resultado típico es timeout en 6817/6818 pese a servicios activos.
\end{criticalnote}

\begin{criticalnote}
\textbf{Uso de perfiles estáticos \texttt{ifcfg} fuera de Ansible.} Editar manualmente \path{/etc/sysconfig/network-scripts/ifcfg-*} en paralelo con este modelo introduce doble fuente de verdad frente a NetworkManager. Los cambios manuales pueden revertirse en la siguiente ejecución del rol o dejar estados híbridos difíciles de depurar.
\end{criticalnote}

\begin{criticalnote}
\textbf{Acoplamiento de grupos de inventario en firewall.} La regla para \texttt{slurmd} se condiciona a \texttt{workers\_r} y \texttt{workers}. Si un compute activo no pertenece a esos grupos, puede quedar sin apertura de 6818 aunque \texttt{slurm\_compute} esté definido en otros plays.
\end{criticalnote}

Un escenario frecuente de degradación parcial es tener interfaces \texttt{int-*} correctamente activas, pero sin rutas persistentes y sin reglas coherentes de firewall. En ese estado, \texttt{ping} o SSH pueden funcionar entre algunos nodos, mientras Slurm falla por timeouts de control. El diagnóstico correcto requiere correlacionar los cinco planos de verificación: interfaz (\texttt{ip a}), perfil NM (\texttt{nmcli}), tabla de rutas (\texttt{ip route}), política efectiva (\texttt{firewall-cmd}) y sockets de servicio (\texttt{ss -lntp}).
